\documentclass[11pt]{article}
\usepackage[a4paper,margin=28mm]{geometry}
\usepackage{amsmath,amssymb,amsthm,mathtools,hyperref}
\newcommand{\ddc}{dd^c}
\newcommand{\PP}{\mathbb P}
\newcommand{\D}{D_\lambda}
\newtheorem{theorem}{Theorem}
\newtheorem{lemma}[theorem]{Lemma}
\newtheorem{proposition}[theorem]{Proposition}
\theoremstyle{remark}\newtheorem{remark}[theorem]{Remark}
\title{Prescribed radial volumes for Gauduchon LCK metrics\newline
on scalar Hopf manifolds}
\author{}
\date{}
\begin{document}
\maketitle
\noindent The unrestricted Gauduchon LCK uniqueness problem is not
resolved by this article.
\begin{abstract}
On each scalar Hopf manifold $H_n$ defined below, $n\geq2$,
a strictly positive smooth radial volume and a negative Lee period
determine exactly one radial Gauduchon locally conformally K\"ahler
metric. A periodic scalar equation gives existence, while a strict
ordering of its solutions gives uniqueness. Every resulting metric has
an explicit positive twisted potential and zero twisted Bott--Chern
class in the actual Lee gauge. A nonconstant radial volume always
produces a non-Vaisman metric. The conclusion concerns the stated
radial family and does not assert uniqueness among all LCK metrics.
\end{abstract}

\section{Forms, class data and statement}
Use $d^c=\frac{i}{4\pi}(\bar\partial-\partial)$ and
$\ddc=\frac{i}{2\pi}\partial\bar\partial$. Put $m=n-1$, where
$n\ge2$, and set
\[
 H_n=(\mathbb C^n\setminus\{0\})/\langle z\mapsto ez\rangle,
 \qquad q:H_n\longrightarrow\PP^m.
\]
Let $\kappa$ be the Fubini--Study form with $\int_{\PP^m}\kappa^m=1$.
The circle-valued function $r=\log\|z\|^2\pmod2$ has a global
differential, and
\[
 \alpha=q^*\kappa=\ddc r,
 \qquad \beta=dr\wedge d^cr.
\]
Here $\alpha$ is closed, semipositive and of rank $m$; $\beta$ is
semipositive of rank one and positive in the complex radial direction.
Their kernels are complementary. Direct differentiation gives
\begin{equation}\label{eq:forms}
 \alpha^n=\beta^2=0,\qquad
 d\beta=-dr\wedge\alpha,\qquad
 d^c\beta=-\alpha\wedge d^cr,\qquad
 \ddc\beta=-\alpha^2.
\end{equation}
On a fiber, $d^cr=d\arg z/(2\pi)$ and $r$ has period two;
therefore $\int_{q^{-1}(x)}\beta=2$ and
$\int_{H_n}\alpha^m\wedge\beta=2$.

Consider precisely the metrics
\begin{equation}\label{eq:radial}
 \omega=C(r)\alpha+D(r)\beta,
 \qquad C,D\in C^\infty(\mathbb R/2\mathbb Z),\quad C,D>0.
\end{equation}
Every such form is LCK, with actual closed Lee form
\begin{equation}\label{eq:lee}
 \lambda=l(r)\,dr,\qquad l=\frac{C'-D}{C},
 \qquad d\omega=\lambda\wedge\omega.
\end{equation}
The Lee form is unique since wedging one-forms with a positive
$(1,1)$ form is injective in complex dimension at least two.
Its period is
\begin{equation}\label{eq:period}
 L=\int_0^2l(r)\,dr=-\int_0^2\frac{D(r)}{C(r)}\,dr<0.
\end{equation}
Thus the Lee form is nonexact. Two radial Lee forms have the same
de Rham class exactly when their periods agree: their difference is
the differential of a smooth periodic function in that case, and the
radial circle detects the period in the converse direction.

For a closed real one-form $\lambda$, use
$d_\lambda=d-\lambda\wedge$,
$d_\lambda^c=d^c-\lambda^c\wedge$ and
$D_\lambda=d_\lambda d_\lambda^c$, where
$\lambda^c=\frac{i}{4\pi}(\lambda^{0,1}-\lambda^{1,0})$.
Locally, $\lambda=df$ gives
$D_\lambda u=e^f\ddc(e^{-f}u)$.
The real twisted Bott--Chern group is the space of smooth real
$d_\lambda$-closed $(1,1)$ forms modulo the image of $D_\lambda$
on smooth real functions. This is the normalized convention for the
weighted group of \cite[Definition 18.34]{OVbook}.

\begin{theorem}\label{thm:main}
For every positive smooth two-periodic function $F$ and every
$L_0<0$, there is exactly one metric of the form \eqref{eq:radial}
which is Gauduchon, has volume
\[
 \omega^n=nF(r)\alpha^m\wedge\beta,
\]
and has Lee period $L_0$. Its actual Lee form admits a positive
smooth global function $U$ such that $\omega=D_\lambda U$.
In particular its twisted Bott--Chern class is zero. If $F$ is
nonconstant, the resulting metric is not Vaisman.
\end{theorem}

Theorem~\ref{thm:main} proves uniqueness within the entire stated
radial family without assuming equality of twisted classes.
It consequently applies to radial candidates satisfying the fixed
class identification of \cite[Question 47.56(3)]{OVbook}.
That question has no radial assumption and remains unanswered here.
The known Vaisman prescribed-volume theorem assumes invariance
under the Lee and anti-Lee fields \cite[Theorem 4.1]{OV24}; the
present scalar argument does not invoke that theorem.

\section{The exact periodic equation}
For any smooth two-periodic scalar $G$, the chain rule and
\eqref{eq:forms} give
\begin{equation}\label{eq:scalar}
 \ddc G=G'\alpha+G''\beta,
 \qquad
 \ddc(G\beta)=-G'\alpha\wedge\beta-G\alpha^2.
\end{equation}
Indeed $d^c(G\beta)=-G\alpha\wedge d^cr$, since
$d^cG\wedge\beta=0$. Taking $d$ proves the second formula with
both displayed signs. Expanding the penultimate and top powers gives
\begin{align*}
 \omega^m&=C^m\alpha^m+mC^{m-1}D\alpha^{m-1}\wedge\beta,\\
 \omega^n&=nC^mD\alpha^m\wedge\beta.
\end{align*}
Consequently
\[
 \ddc(\omega^m)=
 \bigl[(C^m)''-m(C^{m-1}D)'\bigr]\alpha^m\wedge\beta.
\]
The Gauduchon condition is equivalent to
\begin{equation}\label{eq:k}
 C^{m-1}(C'-D)=k
\end{equation}
for a real constant $k$. With $F=C^mD$ fixed, this is the equation
\begin{equation}\label{eq:ode}
 C^mC'=kC+F(r),\qquad D=\frac{F}{C^m}.
\end{equation}
Integration over a period yields
\begin{equation}\label{eq:integral}
 k\int_0^2C\,dr+\int_0^2F\,dr=0,
\end{equation}
so $k<0$. Equations \eqref{eq:lee} and \eqref{eq:k} also give
\begin{equation}\label{eq:L}
 l=kC^{-m},\qquad
 L(C)=-\int_0^2\frac{F(r)}{C(r)^n}\,dr.
\end{equation}

\begin{lemma}\label{lem:order}
For fixed $F$, positive periodic solutions of \eqref{eq:ode}
are unique for a fixed $k<0$. If $k_2>k_1$, their solutions satisfy
$C_2>C_1$ everywhere and $L(C_2)>L(C_1)$.
\end{lemma}
\begin{proof}
At any zero of $R=C_2-C_1$, the equation gives
$R'=(k_2-k_1)C_1^{1-m}>0$. Every zero would thus be a strict
upward crossing. A differentiable periodic function cannot have
only such crossings: their nonzero derivatives make the zeros
isolated, compactness makes their number finite, and a return around
the circle requires a downward crossing. Hence $R$ has constant
strict sign. Equation~\eqref{eq:integral} gives
$\int C_i=-\int F/k_i$ and fixes that sign as positive.
Equation~\eqref{eq:L} now gives the strict inequality of periods.

For equal $k$, initial-value uniqueness for the smooth equation
$C'=kC^{1-m}+FC^{-m}$ on $C>0$ identifies solutions which meet.
Solutions which never meet are strictly ordered, contradicting their
equal integrals in \eqref{eq:integral}. This proves uniqueness.
\end{proof}

\section{Existence and prescribed Lee period}
Fix $k<0$, put $c=-k$, and set $f_0=\min F>0$, $f_1=\max F$.
The equation $C'=C^{-m}(F-cC)$ points inward or tangentially at
both ends of the positive interval
\[
 I_k=[f_0/c,f_1/c].
\]
Its forward flow exists on this compact interval throughout one
period. The time-two map is a continuous self-map of $I_k$;
the intermediate value theorem applied to its difference from the
identity gives a fixed point. This includes the degenerate interval
when $f_0=f_1$. The associated positive solution is two-periodic,
and all derivatives agree at the endpoints by the smooth periodic
equation. Lemma~\ref{lem:order} makes this solution $C_k$ unique.

In particular,
\begin{equation}\label{eq:barriers}
 \frac{f_0}{-k}\le C_k(r)\le\frac{f_1}{-k}.
\end{equation}
The function $L(k)=-\int F/C_k^n$ is strictly increasing by
Lemma~\ref{lem:order}. It is continuous: on a compact neighborhood
of any $k_0<0$, the barriers are uniformly positive and finite;
the equation then bounds all derivatives uniformly. Compactness
gives subsequential limits as $k\to k_0$, and the equation and
uniqueness identify every limit as $C_{k_0}$. Thus $C_k\to C_{k_0}$
uniformly and the integral is continuous. The same barriers imply
\[
 \lim_{k\uparrow0}L(k)=0,
 \qquad \lim_{k\to-\infty}L(k)=-\infty.
\]
Every $L_0<0$ therefore selects exactly one $k$. With
$D=F/C_k^m$, equations \eqref{eq:ode}--\eqref{eq:L} verify every
volume, Lee and Gauduchon requirement in
Theorem~\ref{thm:main}. This is uniqueness among all positive
periodic solutions of the radial equation.

\section{Actual twisted potentials and gauge identification}
For this metric let $l=kC^{-m}<0$ and
$f(r)=\int_0^rl(s)\,ds$ on the real radial cover. Thus
$f(r+2)=f(r)+L_0$. The solution of
\begin{equation}\label{eq:U}
 U'-lU=C
\end{equation}
with
\[
 U(0)=\frac{e^{L_0}}{1-e^{L_0}}
       \int_0^2e^{-f(s)}C(s)\,ds,
 \qquad
 U(r)=e^{f(r)}\left[U(0)+\int_0^re^{-f(s)}C(s)\,ds\right]
\]
is strictly positive on $[0,2]$ and satisfies $U(2)=U(0)$.
Periodicity of the equation makes all derivatives match, so $U$
is a global smooth positive function on $H_n$. Directly, in the
task's exact twisted convention,
\begin{align*}
 D_\lambda U
 &= (U'-lU)\alpha+
    \bigl[U''-2lU'+(l^2-l')U\bigr]\beta\\
 &=C\alpha+(C'-lC)\beta=\omega.
\end{align*}
This verifies the actual zero twisted class without any assertion
about ordinary Bott--Chern exactness.

If $\lambda_2=\lambda_1+da$ has the same radial period,
$a$ is a smooth periodic function, with its additive constant fixed
by the chosen class identification. The exact covariance formula is
\[
 e^{-a}D_{\lambda_2}U_2=D_{\lambda_1}(e^{-a}U_2).
\]
Thus the two identified classes are indeed zero and the class
difference has the genuine potential $e^{-a}U_2-U_1$.
No addition of a constant to a twisted potential is used as a
normalization. For two radial metrics with the same volume and Lee
class, the preceding uniqueness proof already gives equal metrics.

\section{Non-Vaisman volumes and an explicit family}
The squared dual norm of $dr$ for \eqref{eq:radial} is $K/D$,
where $K>0$ is a fixed constant determined by our differential
normalization. Indeed the radial complex line is orthogonal to
$\ker dr\cap\ker d^cr$, and its metric is $D\beta$.
Equation~\eqref{eq:L} therefore gives
\[
 |\lambda|_\omega^2=Kk^2\frac{C^{-2m}}{D}.
\]
A Vaisman Lee form is parallel \cite[Definition 2.3]{OV24},
so its norm is constant. Since $k\ne0$, this would imply
$D=Q C^{-2m}$ for a constant $Q>0$. At both a maximum and a
minimum of $C$, equation~\eqref{eq:k} would then give
\[
 kC^{m+1}+Q=0.
\]
This has only one positive root. The maximum and minimum of $C$
would coincide, forcing $C,D,F$ constant. Every nonconstant $F$
thus gives a non-Vaisman metric, completing the theorem.

For an explicit family take
\[
 C=1+\varepsilon\cos(\pi r),\qquad
 D=C'+C^{1-m},\qquad k=-1,
\]
where
\[
 0<\varepsilon<\min\left\{\tfrac12,
                 \frac{(3/2)^{1-m}}{4\pi}\right\}.
\]
Then $D\ge(3/2)^{1-m}-\pi\varepsilon>
\tfrac34(3/2)^{1-m}>0$. Its exact Lee form and volume density are
\[
 \lambda=-C^{-m}dr,
 \qquad F=C^mC'+C>0.
\]
The Gauduchon identity holds with $k=-1$. Moreover
$F(0)=1+\varepsilon$, $F(1)=1-\varepsilon$ and
$\frac12\int_0^2F\,dr=1$. With
$dV=\frac12\alpha^m\wedge\beta$, a positive probability measure,
the exact volume is $\omega^n=2nF\,dV$.

Its volume is not invariant under the radial Lee translation of
the standard Vaisman metric $\alpha+\beta$, since $F$ is
nonconstant. This distinguishes the prescribed data from the
invariance hypotheses in \cite[Theorem 4.1]{OV24}. It does not
assert noninvariance under the new metric's own Lee field:
its radial volume flux, with the universal factor $4\pi$ omitted,
is $lF/D=-1$, consistently with its Gauduchon condition.
The example is an admissible non-Vaisman
member of the radial uniqueness theorem, not an unequal pair
answering the unrestricted question negatively.

\begin{thebibliography}{9}
\bibitem{OV24} L. Ornea and M. Verbitsky,
\emph{A Calabi--Yau theorem for Vaisman manifolds},
Communications in Analysis and Geometry \textbf{32} (2024),
no. 10, 2889--2900. Primary version:
\href{https://arxiv.org/abs/2206.08808v2}{arXiv:2206.08808v2},
Definition 2.3 and Theorem 4.1 with its proof and footnote 5.
\bibitem{OVbook} L. Ornea and M. Verbitsky,
\emph{Principles of Locally Conformally K\"ahler Geometry},
\href{https://arxiv.org/abs/2208.07188v6}{arXiv:2208.07188v6},
Definition 18.34 and Question 47.56(3), pp. 263 and 683--684.
\end{thebibliography}
\end{document}
